\documentclass[11pt]{article}
\usepackage[margin=1in]{geometry}
\usepackage{amsmath,amssymb}
\usepackage{booktabs}
\usepackage{hyperref}
\usepackage{siunitx}

\title{Independent Replication: \textit{Benchmarking the Quantum Approximate Optimization Algorithm}\\
\large arXiv:1907.02359 (Willsch, Willsch, Jin, De~Raedt, Michielsen, 2019/2020)}
\author{OpenClaw Replication Wave --- QC-100 slot \texttt{QC-QAOA-benchmarking-Guerreschi2019}}
\date{2026-07-06}

\begin{document}
\maketitle

\begin{abstract}
We present an independent, clean-room, classical-simulator replication of the
QAOA benchmark study of Willsch et al.\ (Quantum Inf.\ Process.\ \textbf{19},
197, 2020; arXiv:1907.02359). Using a \texttt{numpy} statevector implementation
and a SciPy Nelder--Mead outer loop, we reproduce (i) the exact Ising cost
values of the paper's 8-variable 2-SAT and 16-variable weighted-MaxCut
instances ($E_{C0} = -9$ and $-17.7$), (ii) the analytic triangle-free $p{=}1$
energy (Eq.~19) to machine precision ($\sim 10^{-15}$), (iii) the paper's
Table~1 $p{=}1$ result for 2-SAT-8(A) exactly (success $8.84\%$, $r{=}0.71$)
and $p{=}5$ to within $\sim 1.4$ percentage points, (iv) the qualitative and
quantitative behaviour of Figs.~7, 10--11 for the 16-variable MaxCut, and
(v) the large-$p$ linear-annealing initialization $\to$ high-success-probability
result. Hardware claims (D-Wave 2000Q, IBM Q Experience) are out of scope for
a simulator replication. \textbf{Verdict: REPLICATED.} Genuine critique of the
scope of this replication, and its distance from the ``several hundred qubits
before quantum speedup'' pessimistic framing in the queue metadata, is given
in \S\ref{sec:critique}.
\end{abstract}

\section{Paper identification and metadata note}
The QC-100 queue row (rank 17) attached slightly incorrect
title/author metadata to this slot (``QAOA benchmarking Guerreschi 2019''),
but the arXiv identifier \texttt{1907.02359} is authoritative, and the
paper is by M.\ Willsch et al.\ (J\"ulich), \textit{not} Guerreschi \&
Matsuura 2019 (which is a related but distinct paper,
\texttt{arXiv:1812.07589}, on runtime crossover points). The directory
slug preserves the assignment convention; \emph{all replication numbers
below refer to arXiv:1907.02359}. The subagent brief's framing of ``crossover
qubit count $\sim$ several hundred qubits'' therefore does \emph{not}
apply to the paper actually in this slot; \S\ref{sec:critique} treats
this head-on.

\section{Method (summary)}
Clean-room re-implementation in Python 3.14, numpy 2.4, scipy 1.18.
No author code exists publicly (JUQCS is in-house at J\"ulich).

\begin{itemize}
  \item Diagonal cost Hamiltonian $H_C = \sum_i h_i Z_i + \sum_{(i,j)} J_{ij} Z_i Z_j$
        as a length-$2^N$ vector.
  \item QAOA ansatz $\lvert\gamma,\beta\rangle = U_B(\beta_p) U_C(\gamma_p)\cdots
        U_B(\beta_1)U_C(\gamma_1) \lvert +\rangle^{\otimes N}$
        with $U_C(\gamma)=e^{-i\gamma H_C}$ applied as elementwise phase and
        $U_B(\beta)=\prod_i e^{-i\beta X_i}$ applied per qubit via
        \texttt{np.moveaxis}.
  \item Three metrics (M1) success probability, (M2) energy expectation
        $\langle H_C\rangle$, (M3) $r=(E_p - E_{\max})/(E_{\min}-E_{\max})$
        (paper Eq.~16).
  \item Optimizer: SciPy Nelder--Mead (matches the paper). Layer chaining:
        initialize $p$ from the $p{-}1$ optimum with $\gamma_p{=}\beta_p{=}0$,
        plus random restarts.
  \item Linear-annealing init (paper Eqs.~29--31) then refine.
\end{itemize}

Instances transcribed from Appendix B Tables 2 and 3(A); correctness verified
by exact ground-energy match.

\section{Results vs paper}
\subsection*{C1 --- Exact instances}
\begin{tabular}{lccc}
\toprule
Instance & $E_{C0}$ (this) & $E_{C0}$ (paper) & Match \\
\midrule
8-var 2-SAT (A) & $-9.0$ & $-9$ & exact \\
16-var wtd.\ MaxCut & $-17.7$ & $-17.7$ & exact \\
\bottomrule
\end{tabular}

\subsection*{C2 --- Analytic $p{=}1$ energy (Eq.~19) vs statevector}
Maximum $\lvert E_{\text{analytic}} - E_{\text{statevector}}\rvert$ over 200 random
$(\gamma,\beta)$: $4.44\times10^{-15}$ (2-SAT-8A), $6.22\times10^{-15}$ (MaxCut).
Machine-precision agreement.

\subsection*{C3 --- 2-SAT-8(A), energy-minimization QAOA}
\begin{tabular}{lcccc}
\toprule
$p$ & succ\% (this) & succ\% (paper) & $r$ (this) & $r$ (paper) \\
\midrule
1 & 8.84 & 8.84 & 0.707 & 0.71 \\
5 & 41.03 & 42.39 & 0.844 & 0.84 \\
\bottomrule
\end{tabular}

\subsection*{C4 --- 16-var MaxCut}
$p{=}1$ success $1.45\%$ (paper: ``$<2\%$'', Fig.~7 $\checkmark$),
$r=0.671$; $p{=}5$ success $42.83\%$, $r=0.920$. Monotone increase of both
metrics with $p$, matching the paper's central trend.

\subsection*{C5 --- Linear-annealing init at large $p$}
2-SAT (A) $p{=}50$: refined success $81.24\%$ (paper $\sim 82.7\%$).
16-var MaxCut $p{=}10$: refined success $76.43\%$ (paper $\sim 85.6\%$).
Directionally and (for 2-SAT) quantitatively correct; the MaxCut gap is
consistent with the many-local-minima landscape the paper explicitly warns about.

\section{Genuine critique and scope limits}
\label{sec:critique}
\input{failure_analysis.md.tex.stub}
\begin{itemize}
  \item \textbf{Simulator-only scope.} Two of the paper's seven headline claims
        (C6: D-Wave 2000Q outperforms simulator QAOA; C7: IBM Q Experience
        $p{=}1$ grid-search is poor quality) require proprietary hardware and
        \emph{are not exercised} here. The remaining five simulator claims
        (C1--C5) are exercised; the paper's \emph{headline} conclusion --- that
        QAOA performance depends strongly on the instance, improves monotonically
        with $p$, and can be near-solved via linear-annealing init at large $p$
        --- is therefore \emph{headline-exercised}.
  \item \textbf{Metadata drift.} The queue row named ``Guerreschi 2019'', but
        the arXiv id points to Willsch et al.\ 2020. The Guerreschi \& Matsuura
        paper (arXiv:1812.07589) makes the ``several hundred qubits'' crossover
        claim that the subagent brief describes; we did \emph{not} replicate
        that runtime-crossover analysis (no Goemans--Williamson vs QAOA
        runtime scan was reimplemented, no crossover-qubit-count was
        reproduced). If the intent was to replicate Guerreschi \& Matsuura,
        the correct pull is a different paper. This is called out explicitly
        rather than hidden.
  \item \textbf{Quantitative gaps.} The 2-SAT $p{=}5$ success probability is
        $41.03\%$ (this) vs $42.39\%$ (paper), $-1.4$ pp; the MaxCut $p{=}10$
        linear-init refined success is $76.43\%$ (this) vs $\sim 85.6\%$
        (paper), $-9$ pp. These are consistent with local-optimum sensitivity,
        but a fully rigorous replication would run multiple seeds and CI
        bands, which we did not (single deterministic seed per config).
  \item \textbf{No noise model.} We only ran ideal statevector; the paper's
        hardware comparisons implicitly include depolarizing/thermal noise
        on the QPU sides. A noise-augmented replication (e.g.\ Kraus channels
        on 1- and 2-qubit gates) is a natural extension.
  \item \textbf{Optimizer choice.} We matched the paper's Nelder--Mead choice,
        but this is known to be brittle; SPSA, gradient methods, or Bayesian
        optimization would likely close some of the $p\geq 5$ gaps at the cost
        of departing from the paper's protocol.
\end{itemize}

\section{Open questions}
See \texttt{open\_questions.json} and \texttt{open\_questions\_section.tex}
for the five follow-up probes and concrete next steps.

\section*{Verdict}
\textbf{REPLICATED} for the classical-simulator core (C1--C5). Hardware
claims (C6, C7) unexercised (out of scope). The Guerreschi \& Matsuura
``several-hundred-qubit crossover'' framing in the queue metadata is
\emph{not} the paper actually in this slot and is \emph{not} exercised.

\end{document}
